\documentclass[11pt]{article}
\input{packages.tex}
\usepackage{ulem}
\DeclareMathOperator{\sign}{sign}
\DeclareMathOperator{\Imm}{Im}
\DeclareMathOperator{\rg}{rg}

% Настройка колонтитулов
\pagestyle{fancy}
\fancyhf{} % Очистить все поля
\fancyhead[L]{Если хочешь победить весь мир, победи \sout{первую пару} себя}
\fancyhead[R]{\href{https://t.me/ceagest}{Автор}}
\fancyfoot[C]{\thepage}               % Справа внизу

\setlength{\headheight}{13.6pt}

\hypersetup{
    colorlinks=true,
    linkcolor=blue,
    citecolor=blue,
    filecolor=blue,
    urlcolor=blue,
}

\begin{document}

\begin{titlepage}
    \centering
    \vspace*{\fill}  % заполняет пространство сверху

    {\Huge Лекция 4. Базис в конечномерных линейных пространствах и его свойства. Матрица перехода. Линейные подпространства. Линейная оболочка. Сумма и пересечение подпространств. Прямая сумма. Формула Грассмана.
    \par}

    \vspace*{\fill}  % заполняет пространство снизу
\end{titlepage}

\setcounter{page}{1}

\section{Базис в конечномерных линейных пространствах и его свойства. Матрица перехода}

\subsection{Базис в конечномерных линейных пространствах и его свойства}

\begin{lemma} \label{lemma:1.1}
    Пусть $\mathcal{L}$ --- конечномерное линейное пространство, $\dim \mathcal{L} = n$. Тогда любой упорядоченный ЛНЗ набор из $n$ векторов является базисом. 
\end{lemma}

\begin{proof}
    Пусть $\{f_1, \ldots, f_n\}$ --- ЛНЗ набор векторов пространства $\mathcal{L}$. Тогда по теореме, 
    доказанной в предыдущей лекции, получаем:
    \[\forall y \in \mathcal{L} \hookrightarrow \{f_1, \ldots, f_n, y\} \text{ --- ЛЗ}\]
    Раз так, то рассмотрим для этого набора некоторую нетривиальную линейную комбинацию, дающую ноль:
    \[\alpha_1 f_1 + \ldots + \alpha_n f_n + \beta y = 0\]
    Заметим, что коэффициент $\beta$ в этом равенстве заведомо не нулевой, так как иначе для набора $\{f_1, \ldots, f_n\}$ существовала бы 
    нетривиальная линейная комбинация, равная нулю, что невозможно, ведь этот набор ЛНЗ. Тогда мы можем перенести в правую часть равенства 
    слагаемое с $y$ и поделить на $\beta \neq 0$, выразив тем самым $y$ через рассматриваемый набор. Но $y \in \mathcal{L}$ был выбран произвольно, 
    откуда по определению базиса получаем требуемое.
\end{proof}

\begin{lemma} \label{lemma:1.2}
    Пусть $\mathcal{L}$ --- конечномерное линейное пространство, $\dim \mathcal{L} = n$, $k < n$. Если $\{f_1, \ldots, f_k\}$ --- ЛНЗ набор, то его 
    можно дополнить до набора, который будет являться базисом в пространстве $\mathcal{L}$.
\end{lemma}

\begin{proof}
    Давайте заметим, что $\exists y \in \mathcal{L}$ такой, что он не представим в виде линейной комбинации векторов $\{f_1, \ldots, f_k\}$.
    Действительно, если бы от противного это было не так, то набор $\{f_1, \ldots, f_k\}$ по определению являлся бы базисом, так как он ЛНЗ и любой вектор пространства по нему раскладывается, что невозможно в силу $k < n$.
    Тогда этот существующий $y$ можно добавить в набор $\{f_1, \ldots, f_k\}$ и он останется ЛНЗ. 
    \\
    Продолжив эту процедуру до тех пор, пока 
    количество векторов в наборе не станет равно $n$, в итоге получим ЛНЗ набор из $n$ векторов в $n$-мерном линейном пространстве $\mathcal{L}$, а значит 
    по предыдущей лемме мы получили базис, что и требовалось.
\end{proof}

\subsection{Матрица перехода}

Пусть в $n$-мерном линейном пространстве $\mathcal{L}$ есть два базиса $\{e_1, \ldots, e_n\}$ и $\{e_1', \ldots, e_n'\}$. Пусть 

\begin{equation*} \label{eq:sec:1.2:mat} \tag{$\ast$}
    \begin{cases}
        e_1' = a_{11} e_1 + \ldots + a_{n1} e_n \\
        \vdots \\
        e_n' = a_{1n} e_1 + \ldots + a_{nn} e_n
    \end{cases}
\end{equation*} \leavevmode
\\
Тогда матрицей перехода от нештрихованного базиса $e$ к штрихованному базису $e'$ будем называть следующую матрицу:

\begin{equation*}
    S = 
    \begin{pmatrix}
        a_{11} & \ldots & a_{1n} \\
        \vdots & \ddots & \vdots \\ 
        a_{n1} & \ldots & a_{nn}
    \end{pmatrix}
\end{equation*} \leavevmode
\\
Как легко видеть, столбцами матрицы перехода $S$ выступают координатные столбцы базисных векторов $e'$ в базисе $e$. Кроме того, условие (\ref{eq:sec:1.2:mat}) 
можно переписать в матричной форме следующим образом:

\begin{equation*}
    \begin{pmatrix}
        e_1' \ldots e_n'
    \end{pmatrix}
    =
    \begin{pmatrix}
        e_1 \ldots e_n
    \end{pmatrix}
    S
    \quad \text{ или } \quad 
    e' = eS
\end{equation*}

\begin{lemma}
    Определитель матрицы перехода отличен от нуля.
\end{lemma}

\begin{proof}
    Так как $e'$ --- базис в линейном пространстве $\mathcal{L}$, то набор $\{e_1', \ldots, e_n'\}$ --- ЛНЗ. Тогда по утверждению 
    из предыдущей лекции координатные столбцы базисных векторов $e'$ в базисе $e$ также ЛНЗ. Значит столбцы матрицы перехода $S$ являются ЛНЗ, 
    откуда $\det S \neq 0$, что и требовалось.
\end{proof}

\begin{lemma}
    Любая невырожденная матрица $S$ соответствующего порядка может являться матрицей перехода 
    между некоторыми двумя базисами в линейном пространстве.
\end{lemma} \leavevmode
\\
Поймём теперь, как связаны координаты векторов в базисах $e$ и $e'$. Пусть вектор $a \in \mathcal{L}$ 
имеет координатные столбцы $x$ и $x'$ в базисах $e$ и $e'$ соответственно. Тогда верно равенство (здесь $e$ и $e'$ --- строки):
\[a = ex = e'x'\]
При этом выше было показано, что $e' = eS$, откуда:
\[ex = e'x' = (eS)x' = e (Sx')\]
Так как координаты вектора в базисе определяются однозначно, то из равенства выше итого получаем:
\[x = Sx'\]
Нетрудно видеть, что это и есть искомая связь координат векторов.

\section{Линейные подпространства. Линейная оболочка}

\begin{definition}
    Пусть $\mathcal{L}$ --- линейное пространство. \textbf{Непустое} подмножество $\mathcal{L}_1 \subseteq \mathcal{L}$ называется 
    \textbf{линейным подпространством} пространства $\mathcal{L}$, если:
    \[\forall x, y \in \mathcal{L}_1 \hookrightarrow x + y \in \mathcal{L}_1 \quad \text{ и } \quad \forall x \in \mathcal{L}_1 \; \; \forall \alpha \in \mathbb{R} \hookrightarrow \alpha x \in \mathcal{L}_1\]
\end{definition}

\begin{lemma}
    Линейное подпространство является линейным пространством.
\end{lemma}

\begin{proof}
    Заметим, что почти все аксиомы линейного пространства тривиально наследуются от всеобъемлющего пространства, а замкнутость 
    следует из определения. Нам нужно показать только существование нейтрального и обратного элементов.
    \\
    Нейтральный элемент $\overline{0}$ содержится в $\mathcal{L}_1$, так как для $\alpha = 0$ по определению линейного подпространства:
    \[\forall x \in \mathcal{L}_1 \hookrightarrow 0 \cdot x = \overline{0} \in \mathcal{L}_1\]
    Обратный элемент существует в силу аналогичных соображений:
    \[\forall x \in \mathcal{L}_1 \hookrightarrow (-1) \cdot x = (-x) \in \mathcal{L}_1\]
\end{proof}

\begin{remark}
    Приведём несколько примеров линейных подпространств:
    \begin{enumerate}
        \item Если всеобъемлющим пространством взять трёхмерное геометрическое пространство, то линейным подпространством в нём будет множество всех векторов, параллельных данной прямой (плоскости).
        \item Само линейное пространство и нулевое пространство всегда являются линейными подпространствами.
        \item Если всеобъемлющим пространством взять множество всех столбов высоты $n$, то линейным подпространством в нём будет множество решений некоторой однородной системы линейных уравнений с $n$ неизвестными.
    \end{enumerate}
\end{remark}

\begin{definition}
    Пусть в линейном пространстве $\mathcal{L}$ задан набор $\{f_1, \ldots, f_m\}$. Множество, состоящее из всевозможных линейных комбинаций этих векторов вида
    \[\alpha_1 f_1 + \ldots + \alpha_m f_m, \quad \alpha_i \in \mathbb{R},\]
    называется \textbf{линейной оболочкой} $\{f_1, \ldots, f_m\}$. Будем обозначать это множество $\left\langle f_1, \ldots, f_m \right\rangle$.
\end{definition}

\begin{remark}
    Множество $\left\langle f_1, \ldots, f_m \right\rangle$ является линейным подпространством пространства $\mathcal{L}$.
\end{remark}

\begin{remark}
    Вновь рассмотрим набор векторов $\{f_1, \ldots, f_m\}$ в линейном пространстве $\mathcal{L}$ и выделим в нём 
    наибольший по размеру ЛНЗ поднабор $\{f_{\alpha_1}, \ldots, f_{\alpha_n}\}$. Тогда этот поднабор 
    является базисом в линейном подпространстве $\left\langle f_1, \ldots, f_m \right\rangle$. Действительно, этот 
    факт следует из определения линейной оболочки и условия максимальности ЛНЗ поднабора.
\end{remark}

\begin{lemma}
    Размерность линейного пространства, представляющего собой линейную оболочку векторов $\{f_1, \ldots, f_m\}$, не 
    превосходит $m$.
\end{lemma}

\begin{proof}
    Требуемое тривиально следует из предыдущего замечания.
\end{proof}

\begin{lemma}
    Пусть $\mathcal{L}$ --- линейное пространство, $\mathcal{L}_1 \subseteq \mathcal{L}$ --- линейное подпространство пространства $\mathcal{L}$, $\dim \mathcal{L}_1 = \dim \mathcal{L}$. Тогда $\mathcal{L}_1 = \mathcal{L}$.
\end{lemma}

\begin{proof}
    Пусть $\dim \mathcal{L}_1 = \dim \mathcal{L} = n$. Пусть $\{e_1, \ldots, e_n\}$ --- базис в $\mathcal{L}_1$. Так как $\mathcal{L}_1 \subseteq \mathcal{L}$, то базис $e$ содержится в пространстве $\mathcal{L}$.
    Тогда по лемме \ref{lemma:1.1} он также будет являться базисом в пространстве $\mathcal{L}$, ведь он ЛНЗ и количество векторов в нём совпадает с размерностью пространства. Значит любой 
    вектор из $\mathcal{L}$ является линейной комбинацией базисных векторов $e$, причём эта линейная комбинация, в свою очередь, обязательно должна лежать в подпространстве $\mathcal{L}_1$ в силу его замкнутости (базис то изначально в $\mathcal{L}_1$ выбирался).
    Это и означает, что $\mathcal{L} \subseteq \mathcal{L}_1$, откуда $\mathcal{L}_1 = \mathcal{L}$, что и требовалось.
\end{proof}

\begin{lemma}
    Пусть $\mathcal{L}$ --- линейное пространство, $\mathcal{L}_1$ --- линейное подпространство пространства $\mathcal{L}$.
    Тогда пространство $\mathcal{L}_1$ можно задать некоторой однородной системой линейных уравнений.
\end{lemma}

\begin{proof}
    Пусть $\{e_1, \ldots, e_k\}$ --- базис в пространстве $\mathcal{L}_1$. По лемме \ref{lemma:1.2} его можно дополнить до базиса в пространстве $\mathcal{L}$ 
    набором $\{e_{k + 1}, \ldots, e_n\}$. Тогда по определению базиса:
    \[\forall x \in \mathcal{L}_1 \hookrightarrow x = \alpha_1 e_1 + \ldots + \alpha_k e_k\]
    \[\forall x \in \mathcal{L} \hookrightarrow x = \alpha_1 e_1 + \ldots + \alpha_k e_k + \alpha_{k + 1} e_{k + 1} + \ldots + \alpha_n e_n\]
    Из этого сразу следует критерий принадлежности некоторого произвольно выбранного вектора из пространства $\mathcal{L}$ подпространству $\mathcal{L}_1$, а именно:
    \[x \in \mathcal{L}_1 \iff \alpha_{k + 1} = \ldots = \alpha_n = 0\]
    Осталось лишь заметить, что на самом деле $\alpha_i = \alpha_i(x) \; \; \forall i \in \{k + 1, \ldots, n\}$, а значит получаем однородную систему линейных уравнений, которая по критерию выше 
    имеет решениями все элементы пространства $\mathcal{L}_1$ и только их:
    \begin{equation*}
        \begin{cases}
            \alpha_{k + 1}(x) = 0 \\
            \vdots \\
            \alpha_n(x) = 0
        \end{cases}
    \end{equation*}
    Легко видеть, что эта система искомая, а значит требуемое доказано.
\end{proof}

\section{Сумма и пересечение подпространств. Прямая сумма}

\begin{definition}
    Пусть $\mathcal{L}_1$ и $\mathcal{L}_2$ --- линейные подпространства линейного пространства $\mathcal{L}$. \textbf{Суммой} подпространств 
    $\mathcal{L}_1$ и $\mathcal{L}_2$ будем называть линейную оболочку $\mathcal{L}_1 \cup \mathcal{L}_2$. Обозначается $\mathcal{L}_1 + \mathcal{L}_2$.
    Все пространства здесь являются конечномерными.
\end{definition}

\begin{remark}
    Под линейной оболочкой множества $\mathcal{L}_1 \cup \mathcal{L}_2$ понимается $\left\langle f_1, \ldots, f_k, g_1, \ldots, g_m \right\rangle$, 
    где $\{f_1, \ldots, f_k\}$ и $\{g_1, \ldots, g_m\}$ --- базисы в пространствах $\mathcal{L}_1$ и $\mathcal{L}_2$ соответственно.
\end{remark}

\begin{lemma}
    Сумма линейных подпространств является линейным подпространством.
\end{lemma}

\begin{proof}
    Требуемое сразу следует из того, что линейная оболочка является линейным подпространством.
\end{proof}

\begin{example}
    Рассмотрим следующие два линейных подпространства $\mathcal{L}_1$ и $\mathcal{L}_2$:
    \begin{figure}[H]
    \centering
    \includegraphics[width=0.54\linewidth]{pictures/Сумма линейных пространств.jpg}
    \end{figure} \leavevmode
    \\
    Легко видеть, что любой вектор плоскости может быть получен как сумма некоторых двух векторов, один из которых 
    принадлежит пространству $\mathcal{L}_1$, а другой --- пространству $\mathcal{L}_2$. Тогда сумма пространств $\mathcal{L}_1$ и $\mathcal{L}_2$ 
    есть вся плоскость, то есть всё объемлющее пространство. 
\end{example}

\begin{remark} \label{rem:3.2}
    Пусть $\{f_1, \ldots, f_k\}$ и $\{g_1, \ldots, g_m\}$ --- базисы в линейных подпространствах $\mathcal{L}_1$ и $\mathcal{L}_2$ пространства $\mathcal{L}$ соответственно.
    Покажем, как найти базис в $\mathcal{L}_1 + \mathcal{L}_2$. По определению суммы линейных подпространств:
    \[\mathcal{L}_1 + \mathcal{L}_2 = \left\langle f_1, \ldots, f_k, g_1, \ldots, g_m \right\rangle\]
    Тогда любой вектор пространства $\mathcal{L}_1 + \mathcal{L}_2$ является линейной комбинацией векторов объединения базисов $f$ и $g$. Но совокупность этих базисов может быть ЛЗ, 
    поэтому нам нужно выбрать из набора $\{f_1, \ldots, f_k, g_1, \ldots, g_m\}$ максимальный по размеру ЛНЗ поднабор, который и будет являться базисом в $\mathcal{L}_1 + \mathcal{L}_2$ просто по определению.
\end{remark}

\begin{definition}
    Пусть $\mathcal{L}_1$ и $\mathcal{L}_2$ --- линейные подпространства линейного пространства $\mathcal{L}$. \textbf{Пересечением} подпространств $\mathcal{L}_1$ и $\mathcal{L}_2$
    будем называть следующее множество:
    \[\mathcal{L}_1 \cap \mathcal{L}_2 := \{x \in \mathcal{L}: x \in \mathcal{L}_1, x \in \mathcal{L}_2\}\]
\end{definition}

\begin{lemma}
    Пересечение линейных подпространств является линейным подпространством.
\end{lemma}

\begin{proof}
    Пусть $\mathcal{L}_1$ и $\mathcal{L}_2$ --- линейные подпространства линейного пространства $\mathcal{L}$. Покажем 
    для начала, что $\mathcal{L}_1 \cap \mathcal{L}_2 \neq \varnothing$. В самом деле:
    \[\overline{0} \in \mathcal{L}_1, \; \overline{0} \in \mathcal{L}_2 \implies \overline{0} \in \mathcal{L}_1 \cap \mathcal{L}_2\]
    Далее увидим, что $\mathcal{L}_1 \cap \mathcal{L}_2$ замкнуто. Покажем для суммы, так как для умножения на число аналогично:
    \begin{equation*}
        x, y \in \mathcal{L}_1 \cap \mathcal{L}_2
        \implies
        \begin{cases}
            x, y \in \mathcal{L}_1 \\
            x, y \in \mathcal{L}_2
        \end{cases}
        \implies 
        \begin{cases}
            x + y \in \mathcal{L}_1 \\
            x + y \in \mathcal{L}_2
        \end{cases}
        \implies
        x + y \in \mathcal{L}_1 \cap \mathcal{L}_2
    \end{equation*}
    Легко видеть, что требуемое доказано.
\end{proof}

\begin{remark}
    Пусть $\mathcal{L}_1$ и $\mathcal{L}_2$ --- линейные подпространства линейного пространства $\mathcal{L}$.
    Покажем, как найти базис в $\mathcal{L}_1 \cap \mathcal{L}_2$. 
    Зададим для этого пространства $\mathcal{L}_1$ и $\mathcal{L}_2$
    однородными системами линейных уравнений, а затем объединим эти системы в одну. Ясно, что полученная система задаёт 
    пространство $\mathcal{L}_1 \cap \mathcal{L}_2$, так как любое её решение является решением как первой, так и второй системы из объединения, а значит 
    принадлежит обоим подпространствам, то есть лежит в их пересечении. Тогда столбцы фундаментальной матрицы этой системы и будут составлять искомый базис в $\mathcal{L}_1 \cap \mathcal{L}_2$ по 
    определению фундаментальной матрицы.
\end{remark}

\begin{lemma}
    Пусть $\mathcal{L}_1, \mathcal{L}_2, \ldots, \mathcal{L}_k$ --- конечномерные линейные подпространства линейного пространства $\mathcal{L}$. Тогда
    \[\dim \left(\mathcal{L}_1 + \mathcal{L}_2 + \ldots + \mathcal{L}_k\right) \leq \dim \mathcal{L}_1 + \dim \mathcal{L}_2 + \ldots + \dim \mathcal{L}_k\]
\end{lemma}

\begin{proof}
    Действительно, требуемое сразу следует из замечания \ref{rem:3.2}, так как сумма размерностей подпространств равна количеству векторов в объединении всех базисов $\mathcal{L}_1, \ldots, \mathcal{L}_k$, а это объединение, в свою очередь, может оказаться ЛЗ набором, а значит 
    не будет являться базисом, откуда и получаем, что размерность суммы уж точно не больше, чем сумма размерностей, что и требовалось.
\end{proof}

\begin{definition}
    Пусть $\mathcal{L}_1, \mathcal{L}_2, \ldots, \mathcal{L}_k$ --- конечномерные линейные подпространства линейного пространства $\mathcal{L}$. Сумма подпространств $\mathcal{L}_1 + \mathcal{L}_2 + \ldots + \mathcal{L}_k$ называется \textbf{прямой}, если 
    \[\dim \left(\mathcal{L}_1 + \mathcal{L}_2 + \ldots + \mathcal{L}_k\right) = \dim \mathcal{L}_1 + \dim \mathcal{L}_2 + \ldots + \dim \mathcal{L}_k\]
    Прямая сумма подпространств $\mathcal{L}_1, \mathcal{L}_2, \ldots, \mathcal{L}_k$ обозначается $\mathcal{L}_1 \oplus \mathcal{L}_2 \oplus \ldots \oplus \mathcal{L}_k$.
\end{definition}

\begin{theorem}[Формула Грассмана]
    Пусть $\mathcal{L}_1$ и $\mathcal{L}_2$ --- конечномерные линейные подпространства линейного пространства $\mathcal{L}$. Тогда
    \[\dim \left(\mathcal{L}_1 + \mathcal{L}_2\right) = \dim \mathcal{L}_1 + \dim \mathcal{L}_2 - \dim \left(\mathcal{L}_1 \cap \mathcal{L}_2\right)\]
\end{theorem}

\begin{proof}
    Будем считать, что подпространства $\mathcal{L}_1$ и $\mathcal{L}_2$ не нульмерные, иначе всё тривиально. Рассмотрим подпространство $\mathcal{L}_1 \cap \mathcal{L}_2$.
    Если оно нульмерно, то равенство очевидно. В самом деле, рассмотрим объединение базисов подпространств $\mathcal{L}_1$ и $\mathcal{L}_2$. Тогда это объединение не может быть ЛЗ, так как иначе 
    вектор одного из базисов мог быть представлен линейной комбинацией остальных, причём коэффициенты при векторах другого базиса были бы не все равны нулю, ведь все векторы одного базиса по определению ЛНЗ в совокупности. Значит, в силу замкнутости подпространств, 
    этот самый разложенный вектор лежал бы и в $\mathcal{L}_1$, и в $\mathcal{L}_2$, что противоречит условию нульмерности пересечения подпространств. Из этих рассуждений следует, что 
    объединение базисов в подпространствах $\mathcal{L}_1$ и $\mathcal{L}_2$ является в этом случае базисом в их сумме, откуда и получаем нужное равенство.
    \\
    Далее полагаем подпространство $\mathcal{L}_1 \cap \mathcal{L}_2$ не нульмерным. Тогда в нём есть базис 
    $\{e_1, \ldots, e_k\}$. Раз так, то давайте, пользуясь леммой \ref{lemma:1.2}, дополним этот базис до базисов в подпространствах $\mathcal{L}_1$ и $\mathcal{L}_2$ наборами 
    $\{f_1, \ldots, f_m\}$ и $\{g_1, \ldots, g_s\}$ соответственно. Мы можем провернуть такую хитрость, так как по определению пересечения подпространств
    \[\mathcal{L}_1 \cap \mathcal{L}_2 \subseteq \mathcal{L}_1 \quad \text{ и } \quad \mathcal{L}_1 \cap \mathcal{L}_2 \subseteq \mathcal{L}_2\]
    Так как мы дополняли до базисов, то наборы $\{e_1, \ldots, e_k, f_1, \ldots, f_m\}$ и $\{g_1, \ldots, g_s, e_1, \ldots, e_k\}$ представляют собой базисы в $\mathcal{L}_1$ и $\mathcal{L}_2$ соответственно.
    \\
    Увидим теперь, что большой набор $\{g_1, \ldots, g_s, e_1, \ldots, e_k, f_1, \ldots, f_m\}$ является ЛНЗ. Для этого рассмотрим некоторую линейную комбинацию векторов этого набора и поймём, на каких наборах значений коэффициентов мы получаем ноль:
    \[\alpha_1 g_1 + \ldots + \alpha_s g_s + \beta_1 e_1 + \ldots + \beta_k e_k + \gamma_1 f_1 + \ldots + \gamma_m f_m = 0\]
    \[\iff\]
    \[\alpha_1 g_1 + \ldots + \alpha_s g_s + \beta_1 e_1 + \ldots + \beta_k e_k = -\left(\gamma_1 f_1 + \ldots + \gamma_m f_m\right)\]
    Пусть $\widetilde{x} = -\left(\gamma_1 f_1 + \ldots + \gamma_m f_m\right)$. В силу замкнутости подпространства $\mathcal{L}_1$ получаем, что $\widetilde{x} \in \mathcal{L}_1$.
    При этом набор $\{g_1, \ldots, g_s, e_1, \ldots, e_k\}$ является базисом в подпространстве $\mathcal{L}_2$, откуда в силу замкнутости уже этого подпространства $\widetilde{x} \in \mathcal{L}_2$.
    Таким образом, $\widetilde{x} \in \mathcal{L}_1 \cap \mathcal{L}_2$, а значит $\alpha_1 = \ldots = \alpha_s = 0$. Действительно, если бы хотя бы один из $\alpha_i$ не обнулялся, то набор $\{g_1, \ldots, g_s, e_1, \ldots, e_k\}$ стал бы ЛЗ (это легко увидеть, рассмотрев разложение $\widetilde{x}$ по базису в подпространстве $\mathcal{L}_1 \cap \mathcal{L}_2$ и приравняв его к уже имеющемуся), что невозможно, так как этот набор является базисом. Рассуждая аналогично, также получаем, 
    что $\gamma_1 = \ldots = \gamma_m = 0$, а значит и $\beta_1 = \ldots = \beta_k = 0$ в силу ЛНЗ базиса. Тогда итого получили, что рассматриваемая линейная комбинация большого набора обнуляется тогда и только тогда, когда эта комбинация тривиальна. Это и означает, что векторы из большого набора являются ЛНЗ.
    \\
    Из полученного выше следует, что базисом в подпространстве $\mathcal{L}_1 + \mathcal{L}_2$ можно взять тот самый набор $\{g_1, \ldots, g_s, e_1, \ldots, e_k, f_1, \ldots, f_m\}$, так как он ЛНЗ и содержит в себе как базис $\mathcal{L}_1$, так и базис $\mathcal{L}_2$, а 
    значит любой вектор линейной оболочки объединения можно представить в виде линейной комбинации векторов этого набора. Тогда, суммируя все наши страдания, получаем:
    \[\dim \left(\mathcal{L}_1 + \mathcal{L}_2\right) = s + k + m, \quad \dim \mathcal{L}_1 = k + m, \quad \dim \mathcal{L}_2 = s + k, \quad \dim \left(\mathcal{L}_1 \cap \mathcal{L}_2\right) = k\]
    Нетрудно видеть, что это победа, товарищи!
\end{proof}


\begin{lemma}[Эквивалентные условия прямой суммы]
    Пусть $\mathcal{L}_1$ и $\mathcal{L}_2$ --- конечномерные линейные подпространства линейного пространства $\mathcal{L}$. Следующие условия эквивалентны:
    \begin{enumerate}
        \item Сумма подпространств $\mathcal{L}_1$ и $\mathcal{L}_2$ является прямой
        \item $\dim\left(\mathcal{L}_1 \cap \mathcal{L}_2\right) = 0$
        \item Базис в $\mathcal{L}_1 + \mathcal{L}_2$ есть объединение базисов $\mathcal{L}_1$ и $\mathcal{L}_2$
        \item $\forall z \in \left(\mathcal{L}_1 + \mathcal{L}_2\right) \hookrightarrow \exists ! x \in \mathcal{L}_1 \; \; \exists ! y \in \mathcal{L}_2: z = x + y$
        \item $\forall x \in \mathcal{L}_1 \; \; \forall y \in \mathcal{L}_2: x, y \neq 0 \hookrightarrow x, y \text{ --- ЛНЗ}$
    \end{enumerate}
\end{lemma}

\begin{proof}
    Требуемое несложно получается из определения прямой суммы и формулы Грассмана, так что смекалистый читатель без труда проведёт доказательство самостоятельно.
\end{proof}

\end{document}